\chapter*{Abstract}
Process mining analyses event logs from information systems in order to understand and enhance real processes. These processes, however, change while an event log captures them. Such concept drift does not only concern the order of activities, but also the working principles of resources and the situation of all running cases together. Existing detectors mostly focus only on the control flow, and publicly available test data does not contain resource drifts.

This thesis introduces an unsupervised, state-based framework that monitors a process through intra-case, resource, and inter-case states. For each perspective, feature extraction, dimensionality reduction, and clustering form the states, and a drift signal compares their distributions over time. Control-flow-based detectors focus on one perspective, whereas the framework detects drifts in all three perspectives, and states describe what changed. Our generator Rheon generates synthetic event logs with injected drifts supporting nine drift types from three perspectives. The library kairo and a dashboard with an AI assistant implement the framework.

We evaluate the framework on 34 logs of the Padella benchmark, with 29 logs where a drift was injected into a real log, and on 45 Rheon logs with 103 drifts. With the hand-tuned configurations, the framework achieves F1-scores from 0.862 to 0.933, and with configurations from the AI agents only from 0.309 to 0.774. The difference is especially caused by precision: Many false positives of the agent configurations are real changes that are unknown to the benchmark or reactions to different perspectives. The states are thereby useful and explainable indicators for concept drifts as long as the appropriate configurations are found. The really difficult problem is, therefore, finding appropriate configurations, and rules of thumb for window sizes and thresholds offer a starting point for that.